\documentclass[options]{philosophersimprint}
%\usepackage{sectsty}
%\linespread{1.2}
\usepackage{xunicode}
\usepackage{xltxtra}
\defaultfontfeatures{Mapping=tex-text}
\setmainfont[
Ligatures={Common}, Numbers={OldStyle}, Variant=01,
BoldFont=LinLibertine_RB.otf,
ItalicFont=LinLibertine_RI.otf,
BoldItalicFont=LinLibertine_RBI.otf
]{LinLibertine_R.otf}
\setmonofont[Scale=0.8]{DejaVuSansMono.ttf}
\usepackage{stmaryrd}
\usepackage[toc,page]{appendix}
\usepackage{amsmath,amsthm,amssymb}
\usepackage[colorlinks=true, citecolor=blue, linkcolor=blue, hyperfootnotes=false]{hyperref}
\usepackage{nameref}
\usepackage{soul}
\usepackage{frame}
\usepackage[style=apa, natbib=true]{biblatex}
\usepackage{scrextend}
\usepackage{mathrsfs}
%\addtokomafont{labelinglabel}{\sffamily}
\makeatletter
\newcommand\labelitem[2]{%
  \item[#1]\def\@currentlabel{#1}\label{#2}}
\makeatother
\usepackage[inline]{enumitem}
\usepackage{framed, csquotes}
\usepackage{tikz}
\usepackage{scalerel}
\usepackage{titlesec}
\newcommand\halo{{{\circ}\mkern-1mu}}
\usepackage{linguex}
\usepackage{parskip}
\usepackage{fancyhdr}
\pagestyle{plain}
% Diagrams
%\usepackage{pgfplots}
\usepackage{tikz}
\usepackage{qtree}




\usepackage{FOL_Yale/notation}

\begin{document}

PHIL 1115: First-order Logic \hfill Yale, Fall 2026\\
Lecture 6, Handout \hfill Dr. Daniel Grimmer

\textbf{1. Why Trees?}

If you've tried your hand at the second problem set, you may be getting annoyed with truth tables. One issue with them is that they are \textit{labor-intensive}. For example, if you're making the truth table for a formula that has $p$, $q$, $r$, and $s$ as atomic formulas, then you have to construct a truth table of 16 lines! If the formula in question contains $n$ different atomic formulas, then we need $2^n$ rows in our truth table. It quickly becomes impractical to construct every such row, even with the help of a computer.

Thus, what we're going to do now is introduce another method for determining whether an argument is valid (or whether a formula is a contradiction, etc.): trees. Typically, making a tree is faster than making a truth table. And building a tree for propositional logic is still a completely mechanical task.

Note: In this lecture, we will define what it means for a formula to be a ``tree''-tautology or a ``tree''-contradiction, etc. These are distinct from the notions of table-tautology and table-contradiction which were introduced in Lecture 4. Whereas the old table method uses the symbol, $\Entails$, the new tree method will use the symbol, $\Proves$.  In the following two lectures we will see that these two methods (tables and trees) are \textit{completely equivalent}: you always get the same answer using one as you would using the other.

\textbf{2. A Motivating Example}

Consider a set $X$ containing three formulas: $p \Conj q$ and $\Neg p \ \Disj \Neg r$ and $\Neg q \Disj s$. Can we make all of them true at the same time? That is, are they \textit{jointly satisfiable}? This would be a 16 line truth table to work out.

Let me show you another method on the board. \textbf{Trees search for ways of simultaneously satisfying a whole list of formulas.} One begins by placing all of the formulas one wants to make true at the top (at the tree's roots, it's upside-down). We then begin to resolve these formulas one-by-one, placing checkmarks next to them.

We want $p \Conj q$ to be true, but this is just the same as demanding that both $p$ and $q$ are true. This means that we can resolve $p \Conj q$ (checking it off) by adding both $p$ and $q$ to our list of requirements. This completes step 1.

\begin{center}
\begin{minipage}[t]{0.12\textwidth}
\centering
\text{Step 1:} \\
\vspace{1em}
\qquad $p \Conj q$\quad\checkmark \\
$\Neg p \ \Disj \Neg r$\qquad \ \\
$\Neg q \Disj s$\qquad \ \\
$\vert$ \\
$p$ \\
$q$
\end{minipage}
\hfill
\begin{minipage}[t]{0.15\textwidth}
\centering
\text{Step 2:} \\
\vspace{1em}
\Tree[.{\qquad$p \Conj q$\quad\checkmark \\
\qquad $\Neg p \ \Disj \Neg r$\quad\checkmark \\
$\Neg q \Disj s$\quad\\
$\vert$ \\
$p$ \\
$q$} [.{$\Neg p$ \\ x} ] [.{$\Neg r$} ] ]
\end{minipage}
\hfill
\begin{minipage}[t]{0.18\textwidth}
\centering
\text{Step 3:} \\
\vspace{1em}
\Tree[.{$p \Conj q$\quad\checkmark \\
$\Neg p \ \Disj \Neg r$\quad\checkmark \\
$\Neg q \Disj s$\quad\checkmark \\
$\vert$ \\
$p$ \\
$q$} [.{$\Neg p$ \\ x} ] [.{$\Neg r$} [.{$\Neg q$ \\ x} ] [.{$s$ \\ o} ] ] ]
\end{minipage}
\end{center}

Next, we want $\Neg p \Disj \Neg r$ to be true, but there are two ways for this to happen. Therefore we have to branch our tree. We can resolve $\Neg p \Disj \Neg r$ (and check it off) by adding one branch which demands $\Neg p$ and one branch that demands $\Neg r$. This completes step 2. But notice what has happened...

The left-most branch of this tree has a contradiction. Reading this branch from top to bottom it demands that $\Neg q \Disj s$ is true as well as $p$ and $q$ and $\Neg p$. But $p$ and $\Neg p$ cannot both be true. This contradiction forces us to close down this branch (marked with an x). The only viable way to continue is via the $\Neg r$ branch.

We can next proceed as before by resolving our final disjunction $\Neg q \Disj s$. There are two ways for this to be true so we have to branch our tree. We can resolve $\Neg q \Disj s$ by adding one branch which demands $\Neg q$ and one branch that demands $s$. Again the left branch closes (marked with an x) due to a contradiction (it demands both $q$ and $\Neg q$). This completes step 3.

But notice that the right branch did not close. Nothing in the right-most branch is contradictory. Moreover, there is nothing left to do with this tree, we have resolved all of its complex formulas. Hence, we can rest assured that no new contradictions will arise to close down this branch. Because the tree is \textit{complete}, we can declare that its last branch is open and mark it with an o.

In sum: \textbf{Trees branch out searching for a way of simultaneously satisfying all of its root formulas.} All branches start their life open. The tree grows and grows by resolving its complex formulas (checking them off). A branch will close (die) if at any point it contains some formula and its negation (neither need be atomic). Once the tree is done growing (once it's complete) we can check if it has any open branches.

\textbf{Satisfiable}: A set of formulas $X$ is jointly \textbf{tree-satisfiable} (denoted $X\notProves$) if and only if some completely resolved tree growing from these roots has an open branch. In the case of a single formula we write $A\notProves$.

Our worked example has shown that the three formulas listed above are jointly satisfiable. We denote the result as $p \Conj q,\quad \Neg p \ \Disj \Neg r,\quad \Neg q \Disj s \notProves$.

\textbf{Contradiction}:\footnote{The fact that our definitions of tree-satisfiability and tree-contradiction are mutually exclusive and exhaustive only holds because of the order-invariance fact discussed below. (See also Restall's ``New notation'' subsection in Ch 4.)} A set of formulas $X$ is jointly a \textbf{tree-contradiction} (denoted $X\Proves$) if and only if some completely resolved tree growing from these roots has all closed branches. In the case of a single formula we write $A\Proves$.

Exercise: Let $A$ be the formula $(p \Conj  q)\Conj  \Neg p$. Use a tree to show that $A$ is a tree-contradiction, $A\Proves$. Intuitively, it should be: it's claiming that both $p$ and $\Neg p$ are true (as well as $q$).

\begin{center}
\begin{minipage}[t]{0.3\textwidth}
This tree only has one branch and it closes. This shows that it is a tree-contradiction; there is no way to make the root formula true.
\vspace{1em}

We can denote this result as $(p \Conj q) \Conj \Neg p\Proves$. This is different from the claim that this formula is a table-contradiction, $(p \Conj q) \Conj \Neg p\Entails$, which we proved in Lecture 4. Which do you prefer: the table method or the tree method?
\end{minipage}
\hfill
\begin{minipage}[t]{0.15\textwidth}
\begin{center}
\vspace{-2em}
\qquad\qquad\qquad\text{Start from $A$:} \\
\Tree[.{\qquad$(p \Conj q) \Conj \Neg p$\qquad\checkmark}
        {\quad$p \Conj q$\quad\checkmark \\ $\Neg p$ \\ $\vert$ \\ $p$ \\ $q$ \\ x} ]
\end{center}
\end{minipage}
\end{center}

\textbf{3. Tree Rules} Let us next take a look at the nine rules we're allowed to use for trees.\footnote{Restall's printed diagrams for negated disjunction and negated conditional are both misprinted, and two of his worked trees carry errors; his prose is sound throughout. The Reading Guide lists them all. Where the book and this table disagree, trust this table.} With a bit of practice, these should all become intuitive (although on exams you will be given a cheat sheet).
\begin{center}
\begin{minipage}[t]{0.18\textwidth}
\textbf{Double Negation}:
\end{minipage}
\hfill
\begin{minipage}[t]{0.27\textwidth}
\begin{center}
\Tree[.{\qquad$\Neg \Neg A$\quad\checkmark} {$A$} ]
\end{center}
\end{minipage}
\end{center}
\vspace{-0.6cm}
\begin{center}
\begin{minipage}[t]{0.18\textwidth}
\textbf{Conjunction and Negated Conjunction}:
\end{minipage}
\hfill
\begin{minipage}[t]{0.27\textwidth}
\begin{center}
\Tree[.{\qquad$A \Conj B$\quad\checkmark} [.{$A$\\ $B$} ] ]
\Tree[.{\qquad$\Neg (A \Conj B)$\quad\checkmark} {$\Neg A$} {$\Neg B$} ]
\end{center}
\end{minipage}
\end{center}
\vspace{-0.6cm}
\begin{center}
\begin{minipage}[t]{0.18\textwidth}
\textbf{Disjunction and Negated Disjunction}:
\end{minipage}
\hfill
\begin{minipage}[t]{0.27\textwidth}
\begin{center}
\Tree[.{\qquad $A \Disj B$\quad\checkmark} {$A$} {$B$} ]
\Tree[.{\qquad$\Neg(A \Disj B)$\quad\checkmark} [.{$\Neg A$\\ $\Neg B$} ] ]
\end{center}
\end{minipage}
\end{center}
\vspace{-0.6cm}
\begin{center}
\begin{minipage}[t]{0.18\textwidth}
\textbf{Conditional and Negated Conditional}:
\end{minipage}
\hfill
\begin{minipage}[t]{0.27\textwidth}
\begin{center}
\Tree[.{\qquad$A \Cond B$\quad\checkmark} {$\Neg A$} {$B$} ]
\Tree[.{\qquad$\Neg(A \Cond B)$\quad\checkmark} [.{$A$\\ $\Neg B$} ] ]
\end{center}
\end{minipage}
\end{center}
\vspace{-0.6cm}
\begin{center}
\begin{minipage}[t]{0.18\textwidth}
\textbf{Biconditional and Negated Biconditional}:
\end{minipage}
\hfill
\begin{minipage}[t]{0.27\textwidth}
\begin{center}
\Tree[.{\qquad$A \Bicond B$\quad\checkmark} {$A$\\$B$} {$\Neg A$ \\ $\Neg B$} ]
\Tree[.{\qquad$\Neg (A \Bicond B)$\quad\checkmark} {$A$\\$\Neg B$} {$\Neg A$ \\ $B$} ]
\end{center}
\end{minipage}
\end{center}

Notice all of the symmetries! We shall explore these on Problem Set 3.

\textbf{4. Making Trees} Here's how to make a tree, in the general case:
\begin{itemize}
    \item Begin by writing a (finite, non-empty) list of formulas. This is called the \textbf{root} of the tree.
    \item A \textbf{branch} in a tree is a path that takes you from the root to a series of nodes, each connected to the one above it. A branch is \textbf{closed} if some formula and its negation, $A$ and $\Neg A$, both appear on it (neither need be atomic). Otherwise it is \textbf{open}. We write `x' beneath every closed branch and `o' beneath every open branch of a complete tree.
    \item You \textbf{resolve} a formula when you break it down using the rules above. Every formula is \textbf{resolvable} except atomic formulas and negations of atomic formulas. After resolving a formula write a check mark next to it, $\checkmark$.
    \item If there is a resolvable but unresolved formula in your tree, then you can expand the tree by resolving the formula and by adding new nodes. In each case, in order to resolve a formula of a given type, you have to extend all of the open branches running through the node of the given formula.
    \item While the size/shape of the resulting tree will depend upon which order you resolve the formulas, the final verdict about whether it has open/closed branches does not. You cannot make a wrong choice.
\end{itemize}

\textbf{Challenge Question:} Suppose that we want to determine whether some set of formulas, $X$, is tree-satisfiable or a tree-contradiction. We build (one of) its trees and note whether or not it has any open branches. But if we had resolved the formulas in a different order, could we have reached the opposite conclusion? I am telling you the answer is ``no'', but you should ask ``why?'' An answer will be given in the following lectures.

\textbf{5. Some Important Definitions}

So far we have seen how trees can tell us if a set of formulas $X$ are jointly tree-satisfiable ($X\notProves$) or are jointly tree-contradictory ($X\Proves$). As we shall now see, we can use these two notions to define corresponding notions of tautology, falsifiability, contingency, and logical equivalence.

\textbf{Tautology}: A formula of propositional logic $A$ is a \textbf{tree-tautology} (denoted $\Proves A$) if and only if its negation, $\Neg A$, is a tree-contradiction (denoted $\Neg A\Proves$). To test if $A$ is a tree-tautology, one develops the tree for $\Neg A$ and checks if all of its branches end up being closed. Importantly, this is a different notion from a table-tautology, which we denoted $\Entails A$.

Exercise: Let $A$ be the formula $p\Cond (p\Disj q)$. Use a tree to show that $A$ is a tree-tautology, $\Proves A$. Intuitively it should be: If $p$ then surely we have either $p$ or $q$ is true (in fact, $p$ is true).

\begin{center}
\begin{minipage}[t]{0.3\textwidth}
There is only one way for a negated conditional, $\Neg(A\Cond B)$, to be true: Its antecedent, $A$, must be true while its consequent, $B$, is false. This gives us $p$ and $\Neg (p \Disj q)$ on the same branch.
\vspace{1em}

Similarly, there is only one way for a negated disjunction, $\Neg(A\Disj B)$, to be true: Both of its disjuncts would need to be false. This gives us $\Neg p$ and $\Neg q$.
\end{minipage}
\hfill
\begin{minipage}[t]{0.15\textwidth}
\begin{center}
\vspace{-2em}
\qquad\qquad\qquad\text{Start from $\Neg A$:}\\
\Tree[.{\qquad$\Neg (p \Cond (p \Disj q))$\qquad\checkmark}
        {$p$ \\ \quad $\Neg (p \Disj q)$\quad\checkmark \\ $\vert$ \\ $\Neg p$ \\ $\Neg q$ \\ x} ]
\end{center}
\end{minipage}
\end{center}

The branch then closes because it contains both $p$ and $\Neg p$. This shows that the root of the tree, $\Neg A$, is a tree-contradiction. But this is just what it means for $A$ to be a tree-tautology. We can denote this result as $\Proves p \Cond (p \Disj q)$. This is different from the claim that this formula is a table-tautology, $\Entails p \Cond (p \Disj q)$, which we proved in Lecture 4. Which method do you prefer: table or tree?

\textbf{Falsifiable:} A formula of propositional logic, $A$, is \textbf{tree-falsifiable} (denoted $\notProves A$) if and only if its negation, $\Neg A$, is tree-satisfiable (denoted $\Neg A\notProves$). This is equivalent to $A$ not being a tree-tautology. To test if $A$ is tree-falsifiable (equivalently, if $\Neg A$ is tree-satisfiable), one develops the tree for $\Neg A$ and checks if one of its branches remains open. Importantly, this is a different notion from table-falsifiability which we denoted $\notEntails A$.

\textbf{Contingent:} A formula $A$ is \textbf{tree-contingent} if and only if it is both tree-satisfiable and tree-falsifiable. To test if $A$ is tree-contingent, one develops the following two trees checking that they both have open branches: an open branch on $A$ confirms that the formula is tree-satisfiable, whereas an open branch on $\Neg A$ confirms that the formula is tree-falsifiable. Importantly, this is a different notion from table-contingency which we defined in Lecture 4.

Exercise: Let $A$ be the formula $p\Conj  (q \Cond p)$. Use \textit{two trees} to show that $A$ is tree-contingent (both tree-satisfiable and tree-falsifiable).

\begin{center}
\begin{minipage}[t]{0.2\textwidth}
\centering
\text{Start from $A$:} \\
\vspace{1em}
\Tree[.{\qquad$p \Conj (q \Cond p)$\quad\checkmark}
        [.{$p$ \\
        \qquad $q \Cond p$\quad\checkmark} {$\Neg q$ \\ o} {$p$ \\ o} ] ]
\end{minipage}
\hfill
\begin{minipage}[t]{0.2\textwidth}
\centering
\text{Start from $\Neg A$:} \\
\vspace{1em}
\Tree[.{\qquad$\Neg (p \Conj (q \Cond p))$\quad\checkmark}
        {$\Neg p$ \\ o} [.{$\Neg (q \Cond p)$\quad\checkmark} {$q$ \\ $\Neg p$ \\ o} ] ]
\end{minipage}
\end{center}

The left tree shows that $A$ is not a contradiction (its open branches show that $A$ can be satisfied). The right tree shows that $A$ is not a tautology (its open branches show that $\Neg A$ can be satisfied, or equivalently $A$ can be falsified).

Compare this proof that $p\Conj  (q \Cond p)$ is tree-contingent with our earlier table-based proof from Lecture 4. Which method do you prefer?

\textbf{Logical Equivalence}: Two formulas $A$ and $B$ are \textbf{tree-equivalent} if and only if $\Neg(A\Bicond B)$ is a tree-contradiction (denoted $\Neg(A\Bicond B)\Proves$). This is equivalent to $A\Bicond B$ being a tree-tautology. Importantly, this is a different notion from table-equivalence which we defined in Lecture 4.

Exercise: Use a tree to show that the formula $p \Cond q$ is tree-equivalent to the formula $\Neg q \Cond \Neg p$. This example is called the \textbf{contrapositive rule}. Compared to the table-based proof from Lecture 4, which method do you prefer?

\begin{center}
\text{Start from $\Neg (A\Bicond B)$:} \\
\vspace{1em}
\Tree[.{\qquad$\Neg ((p \Cond q) \Bicond (\Neg q \Cond \Neg p))$\quad\checkmark}
        [.{$p \Cond q$\quad\checkmark \\ $\Neg (\Neg q \Cond \Neg p)$\quad\checkmark \\ $\vert$ \\ $\Neg q$ \\ $\Neg \Neg p$\quad\checkmark \\ $\vert$ \\ $p$} [.{$\Neg p$ \\ x} ] [.{$q$ \\ x} ] ] [.{$\Neg (p \Cond q)$\quad\checkmark \\ $\Neg q \Cond \Neg p$\quad\checkmark \\ $\vert$ \\ $p$ \\ $\Neg q$} [.{$\Neg \Neg q$\quad\checkmark \\ $\vert$ \\ $q$ \\ x} ] [.{$\Neg p$ \\ x} ] ] ]
\end{center}

Hence, all branches developing out of $\Neg (A\Bicond B)$ close. This means that $A\Bicond B$ is a tree-tautology and hence that $A$ and $B$ are tree-equivalent.

\begin{framed}
\noindent\textbf{Definitions Summary: Tree-Notions and Notation}\\
\begin{tabular}{l | l | l}%{@{}rl@{}}
tree-satisfiable ($X\notProves$) & Open from $X$ & $X\notProves$ \\
tree-contradiction ($X\Proves$) & Closed from $X$ & $X\Proves$\\
tree-tautology ($\Proves A$) & Closed from $\Neg A$ & $\Neg A\Proves$\\
tree-falsifiable ($\notProves A$) & Open from $\Neg A$ & $\Neg A\notProves$\\
tree-contingent & Open $A$ and Open $\Neg A$ & $A\notProves$ and $\Neg A\notProves$\\
tree-equivalent & Closed from $\Neg(A\Bicond B)$ & $\Neg(A\Bicond B)\Proves$\\
\hline
Next Two Lectures: & & \\
tree-invalid ($X\notProves A$) & Open from $X$ and $\Neg A$& $X,\Neg A\notProves$\\
tree-valid ($X\Proves A$) & Closed from $X$ and $\Neg A$ & $X,\Neg A\Proves$
\end{tabular}
\end{framed}

\end{document}